\vspace{-3pt}
\section{Conclusion and Future Work}
\label{sec:conclusion}
\vspace{-5pt}
We introduce \benchmarkt, the first benchmark to thoroughly evaluate LLM-driven retrieval systems for semi-structured knowledge bases (SKBs). Featuring diverse, natural-sounding queries that require context-specific reasoning across diverse domains, \benchmarkh sets a new standard for assessing real-world retrieval systems. We contribute three large-scale retrieval datasets with human-generated queries and an automated pipeline to simulate realistic user queries. Our experiments on \benchmarkh highlight significant challenges for current models in effectively handling textual and relational information. \benchmarkh paves the way for future research to advance complex, multimodal retrieval systems, focusing on reducing retrieval latency and enhancing reasoning abiliites.

Our current SKBs are limited to textual and relational information. Future work should incorporate additional modalities such as images, videos, and speech to provide a more comprehensive information retrieval system. Despite our anonymization efforts, we acknowledge that privacy remains a potential concern when extending this work to other domains with real user data, which should be protected to ensure compliance with privacy regulations.

\newpage